%!TEX root = main.tex

\section{Discussion}
\label{sec:discussion}

Table~\ref{tbl:rules} summarizes the main results of our experiments.
%
It has the form of a set of rules that allows us to choose a good broadcast algorithm for a given deployment scenario.
%
For instance, when there is no mobility, it is better to use \emph{MPR} no matter the density.
%
As our experiments show, in such scenarios \emph{MPR} consumes less energy, guarantees good coverage and uses efficiently the medium.

\begin{table}[ht]
	\footnotesize
	\centering
	
	\begin{tabular}{lr|p{1.9cm}l}
		\textbf{Mobility} & \textbf{Density} &  \textbf{Suggestion} & \textbf{Notes}\\
		
		\midrule
		Static & Any & \emph{MPR} & Fewer control messages \\
		Walking & 5 & \emph{simple flooding} & Best coverage  \\
		Walking & 10  & \emph{MPR} and \emph{ABBA} &   \\
		Walking & 15+ & \emph{CDS-based} &  \\
		
	\end{tabular}
	
	\caption{Rules on which broadcast algorithm to use in different deployment scenarios.}
	\label{tbl:rules}
	% \vspace{-0.3cm}
\end{table}

The rules for deployment scenarios with mobility at walking speed are more complex, as they take into account other external parameters such as the density.
%
To suggest an algorithm, we first aim at increasing coverage, then the energy consumption and we consider that the medium usage is the least important of the metrics.
%
For that reason, we propose \emph{simple flooding} when the density is very low.
%
Finally, the results show that \emph{MPR} and \emph{ABBA} are good alternatives when there is mobility and the density is close to 10.
%
Although for walking speed \emph{ABBA} is outperformed by CDS-based protocols, it is worth considering it as the main alternative for scenarios where nodes move at higher speeds.  
